\section{Norm scales, actual coefficient costs, and low-rank obstructions}\label{sec:barriers}
\subsection{An upper bound from frequency separation}
Fix finite constants $B_\psi,B_\kappa$ such that
\[
0\le\psi_+(v)\le B_\psi(1+|v|)^{-2},\qquad
|\kappa(y)|\le B_\kappa(1+|y|)^{-2}.
\]
\begin{lemma}\label{lem:upper}
For $Q\ge2$,
\[
\norm A\le4B_\psi(H+4Q^2)+3Q^2.
\]
\end{lemma}
\begin{proof}
Distinct reduced fractions $a/q$ are separated on the circle by at least $(4Q^2)^{-1}$, including the gap across zero. The convolution kernel
$\kappa(h/H)\sum_{q,a}\e(ah/q)$ on all integers has Fourier multiplier
\[
m(\theta)=H\sum_{q,a}\sum_{k\in\ZZ}\psi_+\big(H(a/q-\theta-k)\big)\ge0
\]
by Poisson summation for the fixed Schwartz kernel. The points on the right are separated on the real line by at least $s=H/(4Q^2)$. Each half-open unit interval contains at most $1+1/s$ points. Splitting the positive and negative half-lines and summing $(1+j)^{-2}$ gives
$m(\theta)\le4B_\psi H(1+1/s)$.
Compression to $I_x$ cannot increase the norm. Proposition~\ref{prop:decomp} gives $\norm{A+\diag\Delta}\le\norm T$, and
$\max\Delta<\sum_{Q<q\le2Q}q\le3Q^2$. The claimed bound follows.
\end{proof}
Since $H<Q^2$, the physical scales satisfy $\norm A\ll_\psi Q^2$. No cancellation in $\beta,w$ is used.

\subsection{Energy of every actual row}
\begin{proposition}\label{prop:energy}
For the fixed kernel and every sufficiently large $x$, uniformly in $u\in I_x$,
\[
G(u)\asymp_\psi\frac{HQ^2}{\log Q}.
\]
\end{proposition}
\begin{proof}
Continuity and $\kappa(0)=1$ give fixed $0<\rho,b\le1$ with $|\kappa(y)|\ge b$ for $|y|\le\rho$. For $q\nmid u$, put
\[
S_q(u)=\sum_{\substack{t\in I_x\\t\ne u,\ q\nmid t}}|\kappa((u-t)/H)|^2,\quad
C_q(u)=\sum_{\substack{t\in I_x\\t\ne u,\ q\mid u-t}}|\kappa((u-t)/H)|^2.
\]
Squaring gives the exact identity
\[
G(u)=\sum_{\substack{q\in\QQ\\q\nmid u}}a_q^2
\big(S_q(u)+(q-1)(q-3)C_q(u)\big).
\]
Integral comparison yields $S_q(u)\le2B_\kappa^2H/3$ and
$C_q(u)\le2B_\kappa^2H/(3q)$. Since $1+(q-1)(q-3)/q\le q$, the contribution of each prime is at most $(2B_\kappa^2/3)a_q^2qH$.

If $N-1\ge2\rho H$ and $H\ge2Q/\rho$, at least one side of every row contains
$\lfloor\rho H/q\rfloor\ge\rho H/(2q)$ nonzero congruent differences. Their entries are still units, and $(q-2)^2\ge q^2/9$. This prime therefore contributes at least $(\rho b^2/18)a_q^2qH$.

An integer $u\le x=Q^3$ has at most two distinct prime factors greater than $Q$. Removing the primes dividing $u$ from
$S_Q=\sum_{q\in\QQ}a_q^2q$ loses at most $9Q$. The prime number theorem and partial summation give $S_Q\sim(3/2)Q^2/\log Q$. Beyond a common threshold $S_Q\ge18Q$, so
\[
\frac{\rho b^2}{36}HS_Q\le G(u)\le\frac{2B_\kappa^2}{3}HS_Q
\]
for all actual rows.
\end{proof}
The restriction $u\le Q^3$ is essential to this argument. This is not a theorem for equal-length intervals at arbitrary huge origins or for arbitrary coordinate prefixes.

\subsection{A constant-vector witness and the actual coefficients}
\begin{theorem}[Obstruction to coefficient-blind norm bounds]\label{thm:normbarrier}
For the fixed data of Section~\ref{sec:operator},
\[
\norm A\ge(3/2+o(1))\frac{Q^2}{\log Q},\qquad
\norm w^2\sim\frac x2\log x,\qquad
\norm\beta^2\ge\left(\frac12\log\frac{27}{22}+o(1)\right)\frac x{\log x}.
\]
Consequently,
\[
\liminf_{x\to\infty}\frac{\log Q}{xQ^2}\norm A\,\norm\beta\,\norm w
\ge\frac34\sqrt{\log(27/22)}>0.
\]
Any proof using only $B_x\norm\beta\norm w$, with $B_x\ge\norm A$, cannot by itself supply a fixed positive power saving.
\end{theorem}
\begin{proof}
Let $f=N^{-1/2}\1_{I_x}$. Before diagonal deletion, the positive form is
\[
\ip f{(A+\diag\Delta)f}=\frac1N\sum_q q\int\psi_+(v)V_q(v/H)\,dv,
\]
where $V_q(\alpha)$ is the sum of squared deviations of the unit-residue sums
$\sum_{n\equiv r(q),\,n\in I_x}\e(\alpha n)$ from their mean over the units. It is not further divided by $q-1$. Write
$I_x=\{a+1,\ldots,a+N\}$ and $N=mq+s$, $0\le s<q$, and define
\[
\begin{gathered}
G_m(y)=\sum_{k=0}^{m-1}\e(ky),\\
S_j(\alpha)=\e(\alpha(a+j))\bigl(G_m(\alpha q)+\1_{j\le s}\e(\alpha mq)\bigr),
\quad1\le j\le q.
\end{gathered}
\]
The sum $S_j$ is exactly the residue sum for $a+j$, comprising $m$ complete periods and at most one terminal term. If $0<|\alpha|q\le1/2$,
\[
|G_m(\alpha q)|\le(2|\alpha|q)^{-1},\qquad
|\e(\alpha j)-1|\le2\pi|\alpha|q\quad(1\le j\le q).
\]
Each sum thus lies within $\pi+1$ of a common center, and within $1$ when $\alpha=0$. Subtracting the mean minimizes the sum of squared distances, giving $V_q(\alpha)\le q(\pi+1)^2$. The fixed frequency support and the eventual inequality $2Q/H<1/2$ place the whole integral in this range. The positive form is therefore
$O(Q^3/(N\log Q))=O(1/\log Q)$.
On the other hand, the number of unit integers is $N_q=N(1-1/q)+O(1)$, whence
\[
\ip f{Af}=-\sum_{q\in\QQ}(q-2)+O(1/\log Q).
\]
The prime number theorem and partial summation give the operator lower bound.

For $w$, define $c(p)=((p-1)/(p-2))^2-1=O(1/p)$ and extend multiplicatively on squarefree integers. Expand the nonnegative Euler product, discard masks bounded by $1$, and enlarge the interval to obtain
\[
\norm{b_x^{(z)}}^2\le xW_z^2\prod_{p>z}(1+c(p)/p)\ll xW_z^2.
\]
The product converges since $c(p)/p=O(p^{-2})$. Mertens' product theorem gives
$W_z\ll\log z\ll_K\log\log x$. The prime number theorem also gives
$\sum_{u\in I_x}\Lambda(u+2)^2\sim(x/2)\log x$, with proper prime powers contributing a lower-order term. The two-sided triangle inequality for norms proves the asymptotic for $w$, without assuming orthogonality between $b_x^{(z)}$ and $\Lambda$.

For $\beta$, take primes $x^{2/5}<p\le x^{9/20}$ and $x/(2p)<r\le x/p$. Eventually $p<r$, and both exceed $U$. Thus $n=pr\in I_x$, $\Lambda(n)=0$, and the only divisor of $n$ at most $U$ is $1$, so $\beta(n)=-1$. The smaller prime factor is unique, avoiding repeated counts. The ordinary prime number theorem in doubling intervals is uniform over this range. The number of such products is
\[
(1+o(1))\frac x2\sum_{x^{2/5}<p\le x^{9/20}}\frac1{p\log(x/p)}
=\left(\frac12\log\frac{27}{22}+o(1)\right)\frac x{\log x}.
\]
The final equality follows by partial summation and
$\int_{2/5}^{9/20}[a(1-a)]^{-1}\,da=\log(27/22)$.
Multiplying the three lower bounds proves the last assertion.
\end{proof}

\begin{corollary}\label{cor:normalized}
For the specified row-energy normalization,
\[
H^{-1}\ll_\psi\norm Z\ll_\psi(\log Q)/H.
\]
Moreover, $\norm Z\,\norm{D_G^{1/2}\beta}\,\norm{D_G^{1/2}w}$ is still at least
$c_\psi xQ^2/\log Q$ for some $c_\psi>0$.
\end{corollary}
\begin{proof}
Since $A=D_G^{1/2}ZD_G^{1/2}$,
$\norm A/G_{\max}\le\norm Z\le\norm A/G_{\min}$.
Each weighted vector contributes a factor of at least $\sqrt{G_{\min}}$ relative to its unweighted norm. Thus the norm-product expression is at least
\[
(G_{\min}/G_{\max})\norm A\norm\beta\norm w.
\]
Apply Proposition~\ref{prop:energy} and Theorem~\ref{thm:normbarrier}.
\end{proof}
This is a lower bound on an upper-bound expression, not a lower bound on $|\Cx|$. It does not rule out cancellation between the actual coefficients, other preconditioners, or coefficient-sensitive methods.

\subsection{A growing number of negative directions}
\begin{theorem}[A scoped obstruction to low-rank repairs]\label{thm:rank}
Put $\tau_x=Q^2/\log Q$ and
\[
J_x=\left\lfloor N\left(\frac1{4Q}-\frac1H\right)\right\rfloor,\qquad
d_x=2J_x+1,\qquad k_x=d_x-35.
\]
For all sufficiently large $x$, there is a complex subspace $W_x$ such that
\[
\dim_\CC W_x\ge k_x,\qquad
\ip f{Af}\le-\frac{\tau_x}{2}\norm f^2\quad(f\in W_x).
\]
Every matrix $E$ of complex rank less than $k_x$, even if non-Hermitian, satisfies
$\norm{A-E}\ge\tau_x/2$, and $k_x=Q^2/4-x^{11/32}+O(1)$.
\end{theorem}
\begin{proof}
Set $U_0=A+\diag\Delta\ge0$ and take $f_j(n)=N^{-1/2}\e(jn/N)$ for $|j|\le J_x$. These vectors are orthonormal on $N$ consecutive integers, and eventually $d_x<N$. For all $q\le2Q$ and $|v|\le1$, we have $q|j/N+v/H|\le1/2$. The preceding residue-variance bound gives
\[
\ip{f_j}{U_0f_j}\le\frac{C_0}{N}\sum_q q^2,\qquad C_0=(\pi+1)^2.
\]
Let $F$ have these vectors as columns, $B=F^*U_0F\ge0$, and $M_1=\sum_q q$. Then
\[
\tr B\le C_0(1+2Q/N)M_1.
\]
An actual integer is divisible by at most two shell primes, so
\[
\Delta(n)\ge D_*:=\sum_q\frac{q(q-2)}{q-1}-4Q\sim\frac32\tau_x.
\]
The number of eigenvalues of $B$ exceeding $D_*/2$ is at most $2\tr B/D_*<35$ eventually, and eventually $D_*\ge\tau_x$. Map the $[0,D_*/2]$ spectral subspace of $B$ through $F$ to obtain $W_x$. At most $35$ dimensions are lost; on this space the $U_0$ form is at most $D_*/2$, whereas the removed diagonal is at least $D_*$. This proves negativity. If $\rank E<k_x$, then $W_x\cap\ker E$ is nonzero. A unit vector in the intersection gives
$\norm{A-E}\ge|\ip f{(A-E)f}|\ge\tau_x/2$.
Finally, expand $d_x$ using $N=x/2+O(1)$ to obtain the dimension asymptotic.
\end{proof}
The spectral cutoff does not say that the whole low-frequency band is negative, or that repairs above some higher rank must succeed.

Real witnesses require a separate construction; a complex spectral space cannot simply be renamed a real one. In the same symmetric band, form the real orthonormal matrix $V$ from
$f_0$, $(f_j+f_{-j})/\sqrt2$, and $(f_j-f_{-j})/(\ii\sqrt2)$.
The real symmetric matrix $\operatorname{Re}(V^*U_0V)$ is positive semidefinite and has trace $\tr B$. Repeating the spectral cutoff gives a different real subspace of dimension at least $k_x$. In particular, for a real-linear map
\[
E:\RR^{I_x}\longrightarrow\CC^{I_x},\qquad\rank_{\RR}E<k_x,
\]
a unit vector in its kernel intersected with this real negative subspace gives
$\|A-E\|_{\RR^{I_x}\to\CC^{I_x}}\ge\tau_x/2$.
The norm uses the Euclidean unit sphere in the real domain, and rank means real image dimension, not complex rank. The identity
$f^TAf=f^T\operatorname{Re}(A)f$ applies to a quadratic form at one real vector. It cannot replace the mixed scalar $w^TA\beta$ by $w^T\operatorname{Re}(A)\beta$.
